% ===== PRÉAMBULE CANONIQUE — à coller VERBATIM en tête de CHAQUE .tex =====
\documentclass[11pt,a4paper]{article}
\usepackage[utf8]{inputenc}\usepackage[T1]{fontenc}\usepackage{lmodern}
\usepackage[french]{babel}
\usepackage{amsmath,amssymb,mathtools}\usepackage{icomma}
\usepackage{siunitx}\sisetup{locale=FR}  % \qty{}{}, \si{}, \ang{}
\usepackage{xcolor}\usepackage{colortbl}
\usepackage{tikz}\usepackage{pgfplots}\pgfplotsset{compat=1.18}
\usepackage[siunitx,europeanresistors]{circuitikz} % circuits (élec.)
\usepackage[most]{tcolorbox}\usepackage{enumitem}\usepackage{array,booktabs}
\usepackage[a4paper,margin=2.2cm]{geometry}
\usetikzlibrary{arrows.meta,calc,angles,quotes,positioning,patterns,%
  backgrounds,shapes.geometric,decorations.pathmorphing,decorations.markings,intersections}
\definecolor{primary}{RGB}{222,49,99}\definecolor{primarydark}{RGB}{150,22,55}
\definecolor{defcol}{RGB}{0,114,178}\definecolor{methcol}{RGB}{52,92,148}
\definecolor{excol}{RGB}{0,135,90}\definecolor{attncol}{RGB}{200,40,55}
\tcbset{boxbase/.style={enhanced,breakable,boxrule=0.9pt,arc=2.5pt,
  left=3mm,right=3mm,top=2.3mm,bottom=2.3mm,fonttitle=\bfseries,coltitle=white}}
% --- boîtes TUTORIEL (noms EXACTS, ne pas renommer) ---
\newtcolorbox{cle}[1][]{boxbase,colback=defcol!6,colframe=defcol,
  title={\textsc{À retenir}\ifx\relax#1\relax\relax\else\textmd{ : \itshape #1}\fi}}
\newtcolorbox{methode}[1][]{boxbase,colback=methcol!7,colframe=methcol,sharp corners,
  title={\textsc{Méthode}\ifx\relax#1\relax\relax\else\textmd{ --- \itshape #1}\fi}}
\newtcolorbox{exemplebox}[1][]{boxbase,colback=excol!5,colframe=excol,
  title={\textsc{Exemple}\ifx\relax#1\relax\relax\else\textmd{ : \itshape #1}\fi}}
\newtcolorbox{piege}[1][]{boxbase,colback=attncol!6,colframe=attncol,
  title={\textsc{Piège}\ifx\relax#1\relax\relax\else\textmd{ : \itshape #1}\fi}}
\newtcolorbox{remarque}[1][]{boxbase,colback=black!4,colframe=black!45,
  title={\textsc{Remarque}\ifx\relax#1\relax\relax\else\textmd{ : \itshape #1}\fi}}
% --- boîtes EXERCICE / CORRIGÉ (noms EXACTS, auto-comptées) ---
\newtcolorbox[auto counter]{exo}[1][]{boxbase,colback=primary!4,colframe=primary,
  title={\textsc{Exercice}~\thetcbcounter\ifx\relax#1\relax\relax\else\textmd{ --- \itshape #1}\fi}}
\newtcolorbox[auto counter]{corrige}[1][]{boxbase,colback=excol!4,colframe=excol,
  title={\textsc{Corrigé}~\thetcbcounter\ifx\relax#1\relax\relax\else\textmd{ --- \itshape #1}\fi}}
\newcommand{\vv}[1]{\vec{#1}}\newcommand{\ds}{\displaystyle}
\newcommand{\R}{\mathbb{R}}\newcommand{\N}{\mathbb{N}}\newcommand{\Z}{\mathbb{Z}}
\begin{document}
\sisetup{per-mode=symbol}

\begin{corrige}[anatomie optique de l'\oe il]
\begin{enumerate}[label=\alph*)]
  \item une \textbf{lentille convergente} (cornée $+$ cristallin).
  \item la \textbf{rétine} (membrane photosensible au fond de l'\oe il).
  \item l'\textbf{iris} (qui commande l'ouverture de la pupille).
\end{enumerate}
\end{corrige}

\begin{corrige}[vergence de l'\oe il au repos]
\begin{enumerate}[label=\alph*)]
  \item Objet à l'infini $\Rightarrow$ image au foyer image $\Rightarrow f=$ distance $O$--rétine
        $=\qty{16}{\mm}=\qty{0.016}{\m}$.
  \item $C=\dfrac{1}{f}=\dfrac{1}{0,016}=62{,}5~\delta$. La vergence de l'\oe il est très
        grande (focale minuscule).
\end{enumerate}
\end{corrige}

\begin{corrige}[le sens de l'accommodation]
\begin{enumerate}[label=\alph*)]
  \item elle doit \textbf{augmenter} (un objet proche exige plus de vergence pour ramener l'image sur
        la rétine).
  \item le cristallin se \textbf{bombe} (sa courbure augmente, donc $f$ diminue et $C$ augmente).
\end{enumerate}
\end{corrige}

\begin{corrige}[à chaque défaut son verre]
\begin{enumerate}[label=\alph*)]
  \item myopie : lentille \textbf{divergente} ($C<0$).
  \item hypermétropie : lentille \textbf{convergente} ($C>0$).
  \item presbytie : lentille \textbf{convergente} ($C>0$).
\end{enumerate}
\end{corrige}

\begin{corrige}[le myope]
\begin{enumerate}[label=\alph*)]
  \item il voit flou les objets \textbf{lointains} (l'image d'un objet éloigné se forme en avant de la
        rétine) ; il voit bien de près.
  \item un verre \textbf{divergent}.
\end{enumerate}
\end{corrige}

\end{document}
